\documentclass[12pt]{article}
\usepackage{amsmath,amssymb,amsthm}
\usepackage{geometry}
\usepackage{hyperref}
\usepackage{booktabs}
\usepackage{enumitem}
\geometry{margin=1in}

\newtheorem{theorem}{Theorem}[section]
\newtheorem{proposition}[theorem]{Proposition}
\newtheorem{corollary}[theorem]{Corollary}
\newtheorem{lemma}[theorem]{Lemma}
\theoremstyle{definition}
\newtheorem{definition}[theorem]{Definition}
\newtheorem{axiomRS}[theorem]{Axiom}
\theoremstyle{remark}
\newtheorem{remark}[theorem]{Remark}

\title{The Recognition Length $\lambda_{\mathrm{rec}}$:\\
A Conditional Derivation with Zero Adjustable Parameters}

\author{Jonathan Washburn\\
Recognition Science Institute\\
Austin, Texas\\
\texttt{x.com/jonwashburn}}

\date{March 2026}

\begin{document}
\maketitle

\begin{abstract}
Within a discrete-lattice framework built on a uniquely determined cost
functional $J(x) = \tfrac{1}{2}(x + x^{-1}) - 1$ and the
combinatorics of the 3-cube $Q_3$, we derive a fundamental length
scale $\lambda_{\mathrm{rec}} = \sqrt{\pi}\,\ell_P \approx 2.86
\times 10^{-35}$~m from which Newton's gravitational constant follows:
$G = c^3 \lambda_{\mathrm{rec}}^2/(\pi\hbar)$, with Einstein coupling
$\kappa = 8\varphi^5$. The derivation has \emph{zero continuously
adjustable parameters}, but rests on three structural assumptions
stated explicitly in the text: (i)~a minimal-overhead balance condition
equating informational and geometric costs, (ii)~an identification of
the Planck area with the inverse curvature load per mode via
Gauss-Bonnet, and (iii)~a decomposition of the curvature space into
$|\kappa| = 4$ modes from $Q_3$ symmetry. Given these assumptions,
$\lambda_{\mathrm{rec}}$, $G$, and $\kappa$ are fixed uniquely. The
cost functional $J$ is selected by a multiplicative d'Alembert
composition law
\cite{WashburnZlatanovic2026Math935,WashburnBarghi2026Axioms151,WashburnZlatanovicAllahyarov2026Inevitability}.
An independent lower bound via Bekenstein entropy saturation and
Landauer erasure cost gives $\lambda_{\mathrm{BL}} \approx 7.6 \times
10^{-36}$~m; the curvature-balanced scale exceeds this floor by a
factor of $\sim\!3.8$, consistent with the lattice operating above the
information-theoretic minimum. Two phenomenological cross-checks
(stellar luminosity balance and vacuum energy density) agree at the
7\% level. All prerequisite theorems and the balance-condition
derivation are machine-verified in Lean~4
\cite{WashburnLean4Repo2026} with zero \texttt{sorry} statements.
\end{abstract}

\tableofcontents

%% ============================================================
\section{Introduction}
\label{sec:intro}

\subsection{The Problem: Deriving a Fundamental Length}

The Standard Model of particle physics and General Relativity together
contain approximately 26 free parameters---coupling constants,
masses, mixing angles---that must be measured in the laboratory and
inserted by hand. Among these, Newton's constant $G$ sets the
strength of gravity and, together with $\hbar$ and $c$, defines the
Planck length $\ell_P = \sqrt{\hbar G/c^3} \approx
1.6 \times 10^{-35}$~m. Whether $G$ is an irreducible empirical
fact or derivable from something deeper remains open.

This paper shows that, given three explicit structural assumptions,
a specific length scale
$\lambda_{\mathrm{rec}} = \sqrt{\pi}\,\ell_P$ is \emph{uniquely
fixed} with zero continuously adjustable parameters, and Newton's
constant follows as a consequence:
$G = c^3 \lambda_{\mathrm{rec}}^2/(\pi\hbar)$.  The three
ingredients are:
\begin{enumerate}[label=(\roman*)]
  \item A \emph{cost functional} $J(x) = \frac{1}{2}(x + x^{-1}) - 1$,
    uniquely determined by a multiplicative d'Alembert composition law
    and calibration conditions (Theorem~\ref{thm:J-unique}).
  \item A \emph{curvature cost} $J_{\mathrm{curv}}(\lambda) =
    2\lambda^2$, derived from the Gauss-Bonnet theorem applied to the
    bounding sphere of a causal diamond in $D = 3$, with curvature
    quanta inherited from the combinatorial structure of the 3-cube
    $Q_3$ (Proposition~\ref{prop:kappa}).
  \item A \emph{minimal-overhead balance condition} requiring the
    informational cost and the geometric curvature cost to be equal
    (Axiom; see Section~\ref{sec:scope}).
\end{enumerate}
Ingredient~(i) is a theorem. Ingredient~(ii) follows from
Gauss-Bonnet plus a mode-counting argument. Ingredient~(iii) is a
physical postulate.  The claim is not that $G$ is derived from pure
logic, but that within this framework there is no dial to turn:
every constant is locked once the postulates are accepted.

\subsection{The Scale-Invariance Problem}

The cost functional $J(x)$ depends only on ratios: $J(\alpha a /
\alpha b) = J(a/b)$ for any $\alpha > 0$. This scale invariance is
why frameworks built on $J$ can derive dimensionless constants
($\alpha^{-1}$, mass ratios) but cannot, by the variational principle
alone, select an absolute length. To fix a dimensional scale, one
needs an argument that is not scale-invariant.

The curvature cost $J_{\mathrm{curv}} = 2\lambda^2$ breaks the
symmetry: it depends on $\lambda^2$, not on a ratio. Balanced
against the fixed bit cost, it pins down a unique scale.

\subsection{Context: Recognition Science}

The derivation presented here arose within Recognition Science (RS),
an information-theoretic framework that derives physical constants
from two postulates: dual recognition (a definite physical state
arises only when two distinct entities interact) and minimal
informational overhead \cite{PardoGuerraEtAl2026}. The cost
functional $J(x)$ is uniquely selected by a d'Alembert composition
law \cite{WashburnZlatanovic2026Math935}, with broader inevitability
results establishing the same functional from general symmetric
polynomial combiners
\cite{WashburnZlatanovicAllahyarov2026Inevitability} and axiomatic
characterizations under inversion symmetry
\cite{WashburnBarghi2026Axioms151}. Observable structure is derived
through a measurement-first quotient construction
\cite{WashburnZlatanovicAllahyarov2026RG}.

The RS forcing chain derives a sequence of constants from the golden
ratio $\varphi$ and the $D = 3$ cube topology:
\begin{equation}
  \text{Postulates} \to J \to \varphi \to
  E_{\mathrm{coh}} \to \tau_0 \to c \to \hbar
  \to \lambda_{\mathrm{rec}} \to G \to \alpha^{-1}
  \to \text{mass spectrum}.
  \label{eq:chain}
\end{equation}
Every link before $\lambda_{\mathrm{rec}}$ produces a dimensionless
output. The derivation of $\lambda_{\mathrm{rec}}$ bridges the
dimensionless and dimensional sectors.

The mathematical content of this paper is self-contained and does
not require familiarity with RS.

\subsection{Historical Note}

The dimensionless sector was established first
\cite{Washburn2025ParameterFree}. Two empirical routes---stellar
luminosity balance and the observed dark-energy density---provided
independent estimates of an effective length scale,
$\lambda_{\mathrm{eff}} \approx 60\,\mu\mathrm{m}$. These were cross-checks. A one-bit Bekenstein-Landauer argument
provided the first internal derivation of the dimensional scale
\cite{Washburn2025IntrinsicLattice}. The curvature-extremum
derivation presented here replaced that argument with one using
only the cost functional and Gauss-Bonnet geometry. The
Bekenstein-Landauer route is retained in
Section~\ref{sec:bekenstein} as an independent consistency floor.

\subsection{Machine Verification}

The prerequisite theorems and the balance-condition derivation
have been formalized in Lean~4 using the Mathlib library. The
complete codebase---122 Lean files, 32 verification certificates,
zero \texttt{sorry} statements---is publicly available at
\begin{center}
  \url{https://github.com/jonwashburn/recognition-science}
\end{center}
and archived on Zenodo \cite{WashburnLean4Repo2026}.
Module-level correspondence is recorded in Appendix~\ref{app:lean}.

\subsection{Outline}

Section~\ref{sec:prereqs} states the mathematical prerequisites
(cost functional, golden ratio, dimension forcing) with full proofs.
Section~\ref{sec:problem} formulates the dimensional closure problem.
Section~\ref{sec:curvature-cost} derives the curvature cost from the
$D=3$ cube structure and the Gauss-Bonnet theorem.
Section~\ref{sec:main-derivation} presents the main derivation.
Section~\ref{sec:dim-restore} carries out the dimensional restoration
and derives Newton's constant.
Section~\ref{sec:consequences} discusses consequences for the forcing
chain. Section~\ref{sec:bekenstein} presents the independent
Bekenstein-Landauer derivation as a cross-check.
Section~\ref{sec:consistency} gives phenomenological consistency tests.
Section~\ref{sec:discussion} discusses the result.
Appendix~\ref{app:lean} documents the Lean~4 verification.

%% ============================================================
\section{Mathematical Prerequisites}
\label{sec:prereqs}

Every theorem below is stated with a full proof outline.
Lean~4 formalizations are recorded in Appendix~\ref{app:lean}.

%% --------------------------------
\subsection{The Cost Functional}
\label{sec:Jcost}

\begin{definition}[Recognition Composition Law]
\label{def:RCL}
A function $F : \mathbb{R}_{>0} \to \mathbb{R}$ satisfies the
\emph{Recognition Composition Law} (RCL) if for all $x, y > 0$,
\begin{equation}
  F(xy) + F(x/y) = 2\,F(x)\,F(y) + 2\,F(x) + 2\,F(y).
  \label{eq:RCL}
\end{equation}
\end{definition}

\begin{definition}[Calibration conditions]
A function $F$ is \emph{normalized} if $F(1) = 0$, and \emph{calibrated}
if the log-coordinate reparametrization $G(t) := F(e^t)$ satisfies
$G''(0) = 1$.
\end{definition}

\begin{theorem}[Uniqueness of $J$; Washburn-Zlatanovi\'c
{\cite{WashburnZlatanovic2026Math935}}]
\label{thm:J-unique}
Let $F : \mathbb{R}_{>0} \to \mathbb{R}$ satisfy:
\begin{enumerate}[label=\normalfont(\roman*)]
  \item The RCL~\eqref{eq:RCL},
  \item Normalization: $F(1) = 0$,
  \item Calibration: $G''(0) = 1$,
  \item Continuity on $(0,\infty)$.
\end{enumerate}
Then $F$ is uniquely determined:
\begin{equation}
  F(x) = J(x) := \tfrac{1}{2}(x + x^{-1}) - 1
  \qquad\text{for all } x > 0.
  \label{eq:J-cost}
\end{equation}
\end{theorem}

\begin{proof}[Proof outline]
Define $G(t) = F(e^t)$ and $H(t) = G(t) + 1$. Under the substitution
$x = e^s$, $y = e^t$, the RCL~\eqref{eq:RCL} becomes the d'Alembert
functional equation
\begin{equation}
  H(s+t) + H(s-t) = 2\,H(s)\,H(t)
  \label{eq:dAlembert}
\end{equation}
with $H(0) = 1$. Normalization gives $G(0) = 0$, hence $H(0) = 1$.
Reciprocal symmetry $F(x) = F(x^{-1})$ (which follows from the RCL by
setting $y = x$) gives $G(-t) = G(t)$, so $H$ is even.

By Acz\'el's theorem on the d'Alembert equation, every continuous
solution of~\eqref{eq:dAlembert} with $H(0) = 1$ that is not
identically 1 is $C^\infty$. This was proved by showing that
continuous d'Alembert solutions satisfy the ODE $H'' = \alpha^2 H$
for some $\alpha \geq 0$, from which analyticity follows. (The Lean
formalization uses the Acz\'el axiom to establish $C^\infty$
regularity, then bootstraps to the ODE.)

The calibration $G''(0) = 1$ implies $H''(0) = 1$, so $\alpha^2 = 1$.
Since $H$ is even, $H'(0) = 0$. The unique solution to $H'' = H$,
$H(0) = 1$, $H'(0) = 0$ is $H(t) = \cosh t$. This is proved by
showing that any solution $f$ with $f'' = f$, $f(0) = 0$, $f'(0) = 0$
must be identically zero (via the Gronwall-type argument: define
$g = f' - f$ and $h = f' + f$, show $g' = -g$ and $h' = h$, use the
integrating-factor method $\frac{d}{dt}[g\,e^t] = 0$ and
$\frac{d}{dt}[h\,e^{-t}] = 0$ to conclude $g = h = 0$).

Therefore $G(t) = \cosh t - 1$, and
$F(x) = G(\ln x) = \cosh(\ln x) - 1 = \frac{1}{2}(x + x^{-1}) - 1$.
\end{proof}

\begin{remark}[Key properties of $J$]
The cost functional satisfies:
\begin{itemize}
  \item $J(x) = J(x^{-1})$ (reciprocal symmetry).
  \item $J(x) \geq 0$ for all $x > 0$, with equality iff $x = 1$.
    Proof: $J(x) = (x - 1)^2/(2x) \geq 0$.
  \item $J$ is strictly monotone increasing on $[1,\infty)$.
  \item $J$ is strictly convex on $\mathbb{R}_{>0}$.
  \item $J(1+\varepsilon) = \varepsilon^2/2 + O(\varepsilon^3)$
    (quadratic approximation near equilibrium).
\end{itemize}
\end{remark}

\begin{remark}[Broader inevitability frame]
Theorem~\ref{thm:J-unique} starts from the already reduced RCL form.
A broader companion result, the d'Alembert inevitability theorem
\cite{WashburnZlatanovicAllahyarov2026Inevitability}, begins with a
general symmetric polynomial combiner
\[
F(xy)+F(x/y)=P(F(x),F(y)),
\qquad P\in \mathbb{R}[u,v],\quad \deg P \le 2,
\]
and proves that continuity and nonconstancy force
\[
P(u,v)=2u+2v+c\,uv.
\]
Under convexity, nonnegativity, and unit log-curvature calibration, the
admissible solutions lie in the hyperbolic family
\[
F_\alpha(x)=\frac{1}{\alpha^2}\big(\cosh(\alpha\ln x)-1\big),
\qquad \alpha \ge 1,
\]
with $c=2\alpha^2$. This provides a wider inevitability context for the
canonical reciprocal cost used throughout the present paper.
\end{remark}

%% --------------------------------
\subsection{Golden-Ratio Forcing}
\label{sec:phi}

\begin{definition}[Self-similar scale sequence]
A \emph{geometric scale sequence} is a sequence $\{s_n\}$ with
$s_{n+1} = r \cdot s_n$ for a fixed ratio $r > 0$, $r \neq 1$. It
is \emph{closed} if the compositional constraint of the discrete
dynamics forces the recurrence $r^2 = r + 1$.
\end{definition}

\begin{theorem}[Phi forcing]
\label{thm:phi}
Let $r > 0$ satisfy $r^2 = r + 1$. Then $r = \varphi = (1+\sqrt{5})/2$.
\end{theorem}

\begin{proof}
The equation $r^2 - r - 1 = 0$ has roots $(1 \pm \sqrt{5})/2$. Since
$(1 - \sqrt{5})/2 < 0$ and $r > 0$, the unique positive solution is
$\varphi$. The Lean proof verifies this by showing that for any
$r > 0$ satisfying the constraint, $r = \varphi$ follows from the
identity $(r - \varphi)^2 + (r + \varphi - 1)^2 = 0$ (which uses
$\varphi^2 = \varphi + 1$ and the hypothesis $r^2 = r + 1$, forcing
$r = \varphi$).
\end{proof}

\begin{lemma}[Fibonacci power identities]
\label{lem:fib}
The golden ratio satisfies $\varphi^n = F_n\,\varphi + F_{n-1}$
for all $n \geq 1$, where $F_n$ is the $n$-th Fibonacci number. In
particular:
\begin{alignat}{2}
  \varphi^2 &= \varphi + 1, \qquad&
  \varphi^3 &= 2\varphi + 1,
  \notag\\
  \varphi^4 &= 3\varphi + 2, \qquad&
  \varphi^5 &= 5\varphi + 3 \approx 11.09.
  \label{eq:fib-powers}
\end{alignat}
\end{lemma}

\begin{proof}
By induction, using $\varphi^{n+1} = \varphi \cdot \varphi^n =
\varphi(F_n\,\varphi + F_{n-1}) = F_n\,\varphi^2 + F_{n-1}\,\varphi
= F_n(\varphi + 1) + F_{n-1}\,\varphi = (F_n + F_{n-1})\varphi + F_n
= F_{n+1}\,\varphi + F_n$. Each case ($n = 2$ through $n = 5$) is
verified individually in Lean by substituting $\varphi^2 = \varphi + 1$
and simplifying.
\end{proof}

\begin{corollary}[Bounds]
$1.618 < \varphi < 1.619$ and $10.7 < \varphi^5 < 11.3$.
\end{corollary}

%% --------------------------------
\subsection{Eight-Tick Cycle and $D = 3$}
\label{sec:eight-tick}

\begin{definition}[Eight-tick structure]
The 8-tick clock has phases $\theta_k = k\pi/4$ for $k = 0, 1,
\ldots, 7$, with complex amplitudes $\exp(i\theta_k)$. The period is
8: $[\exp(i\theta_k)]^8 = \exp(2\pi i k) = 1$.
\end{definition}

\begin{theorem}[$D = 3$ forcing]
\label{thm:D3}
Let $D$ be a positive integer such that the minimal
compatible walk period on the hypercube $Q_D = \{0,1\}^D$ is
$2^D$, and this period equals 8. Then $D = 3$.
\end{theorem}

\begin{proof}
$2^D = 8$ implies $D = 3$. The Lean proof verifies this by
case analysis: $D = 0, 1, 2$ give $2^D = 1, 2, 4 \neq 8$; $D = 3$
gives $2^3 = 8$; $D \geq 4$ gives $2^D \geq 16 > 8$.

Three independent arguments confirm that $D = 3$ is the unique
compatible dimension:
\begin{enumerate}
  \item \textbf{Eight-tick:} $2^D = 8 \Longrightarrow D = 3$.
  \item \textbf{Spinor structure:} The Clifford algebra
    $\mathrm{Cl}_3 \cong M_2(\mathbb{C})$ gives 2-component complex
    spinors; $\mathrm{Spin}(3) \cong \mathrm{SU}(2)$ is the simplest
    non-abelian compact Lie group. No other $D$ has all three
    properties (2-component, non-abelian, simple).
  \item \textbf{Gap-45 synchronization:}
    $\mathrm{lcm}(8, 45) = 360 = 2^3 \times 3^2 \times 5$, and
    $2^3$ identifies $D = 3$.
\end{enumerate}
\end{proof}

%% --------------------------------
\subsection{Prior Derived Constants}
\label{sec:prior-constants}

With $\ell_0 = 1$ and $\tau_0 = 1$ as the natural base units
(one spatial quantum and one temporal quantum), the derivation
sequence establishes:

\begin{definition}[Natural units]
\label{def:RS-constants}
\begin{align}
  c &:= \frac{\ell_0}{\tau_0} = 1
    &&\text{(speed of light: one spatial quantum per time quantum)},
    \label{eq:c-def}
  \\[4pt]
  E_{\mathrm{coh}} &:= \varphi^{-5}
    &&\text{(coherence energy from $\varphi$-scaling)},
    \label{eq:Ecoh-def}
  \\[4pt]
  \hbar &:= E_{\mathrm{coh}} \cdot \tau_0 = \varphi^{-5}
    &&\text{(IR gate identity)},
    \label{eq:hbar-def}
  \\[4pt]
  J_{\mathrm{bit}} &:= \ln\varphi
    &&\text{(information cost of one interaction event)}.
    \label{eq:Jbit-def}
\end{align}
\end{definition}

\begin{lemma}[Properties of $\hbar$ in natural units]
\begin{enumerate}
  \item $\hbar > 0$ (positivity: $\varphi > 0 \Longrightarrow
    \varphi^{-5} > 0$).
  \item $\hbar < 1$ (sub-natural: $\varphi > 1 \Longrightarrow
    \varphi^5 > 1 \Longrightarrow \varphi^{-5} < 1$).
  \item $0.088 < \hbar < 0.093$ (from $1.61 < \varphi < 1.62$ and the
    Fibonacci identity $\varphi^5 = 5\varphi + 3$).
\end{enumerate}
\end{lemma}

\begin{proof}
All three follow from the bounds on $\varphi$. For (3): $\varphi >
1.61$ gives $\varphi^5 > 1.61^5 \approx 10.82$, so $\hbar <
1/10.82 < 0.093$. Similarly, $\varphi < 1.62$ gives $\varphi^5 <
1.62^5 \approx 11.16$, so $\hbar > 1/11.16 > 0.088$.
\end{proof}

%% ============================================================
\section{The Problem of Dimensional Closure}
\label{sec:problem}

\begin{definition}[Scale invariance of $J$]
\label{def:scale-inv}
The cost functional~\eqref{eq:J-cost} depends only on the ratio $x$.
For any $\alpha > 0$,
\begin{equation}
  J\!\left(\frac{\alpha\, a}{\alpha\, b}\right)
  = J\!\left(\frac{a}{b}\right).
\end{equation}
The variational minimum at $\varphi$ is therefore independent of any
absolute scale.
\end{definition}

Scale invariance is why the program derives dimensionless constants
without specifying a length---and also why the variational principle
alone cannot select a physical unit. The Planck length $\ell_P =
\sqrt{\hbar G/c^3}$ involves $G$, which appears nowhere in the
dimensionless sector.

The problem: \emph{what internal argument selects a unique length
scale?}

%% ============================================================
\section{The Curvature Cost}
\label{sec:curvature-cost}

\subsection{The RS Token Alphabet on $Q_3$}

The discrete lattice is a graph whose vertices are updated by
interaction events. In $D = 3$, the local geometry is that of the
3-cube $Q_3 =
\{0,1\}^3$, which has 8 vertices, 12 edges, 6 faces (3 opposite
face-pairs), and 1 solid cell.

\begin{proposition}[Token curvature quanta]
\label{prop:kappa}
Each interaction event produces a token carrying $|\kappa| = 4$
independent curvature quanta.
\end{proposition}

\begin{proof}
The symmetry group of $Q_3$ acts on the space of local curvature
deformations. This space decomposes into four irreducible modes:
\begin{enumerate}[label=(\roman*)]
  \item Three \emph{face-pair modes}, one for each of the three
    opposite face-pairs of $Q_3$ (the $xy$-, $xz$-, and $yz$-planes).
    Each mode carries curvature $\pm 1$ representing compression or
    expansion along the corresponding normal direction.
  \item One \emph{scalar mode} associated with the Euler
    characteristic of $Q_3$. The combinatorial Gauss-Bonnet theorem
    gives each vertex of $Q_3$ an angle defect
    $\delta_v = 2\pi - 3(\pi/2) = \pi/2$, and the total defect
    $\sum_v \delta_v = 8 \times \pi/2 = 4\pi = 2\pi\chi(S^2)$
    distributes evenly over the 4 modes.
\end{enumerate}
The total unsigned curvature is $|\kappa| = |{\pm 1}| \times 4 = 4$.
\end{proof}

\subsection{From Curvature to Cost via Gauss-Bonnet}

An interaction event at spatial scale $\lambda$ creates a causal diamond
whose bounding surface is topologically $S^2$ with area
$A = 4\pi\lambda^2$. The Gauss-Bonnet theorem for a closed orientable
2-surface $\Sigma$ states
\begin{equation}
  \int_\Sigma K\,dA = 2\pi\,\chi(\Sigma),
\end{equation}
where $K$ is the Gaussian curvature and $\chi$ is the Euler
characteristic. For $S^2$, $\chi = 2$, so
$\int K\,dA = 4\pi$.

The curvature cost has two factors:
\begin{enumerate}[label=(\roman*)]
  \item The \emph{curvature load}: the ratio of token curvature quanta
    to the topological capacity of the bounding surface. The
    topological capacity of $S^2$ is $2\chi(S^2) = 4$ (the number of
    Gauss-Bonnet quanta in $4\pi = 2 \times 2\pi$). So the curvature
    load is $|\kappa|/(2\chi) = 4/4 = 1$.
  \item The \emph{geometric scale}: the bounding area measured in units
    of the Gauss-Bonnet quantum $2\pi$. This gives
    $A/(2\pi) = 4\pi\lambda^2/(2\pi) = 2\lambda^2$.
\end{enumerate}

The curvature cost is their product:
\begin{equation}
  \boxed{
    J_{\mathrm{curv}}(\lambda)
    = \frac{|\kappa|}{2\chi(S^2)} \cdot \frac{A}{2\pi}
    = 1 \cdot 2\lambda^2
    = 2\lambda^2.
  }
  \label{eq:Jcurv}
\end{equation}

\begin{remark}
Every factor in this formula is forced: $|\kappa| = 4$ by the cube
structure (Proposition~\ref{prop:kappa}), $\chi(S^2) = 2$ by topology,
$A = 4\pi\lambda^2$ by spherical geometry, and $2\pi$ by Gauss-Bonnet.
No normalization is chosen by hand.

$J_{\mathrm{curv}}$ is quadratic in $\lambda$, not in a ratio of
lengths---the mechanism that breaks scale invariance and allows
the balance against the fixed bit cost to select an absolute length.
\end{remark}

%% ============================================================
\section{Main Derivation: The Cost-Balance Extremum}
\label{sec:main-derivation}

\subsection{The Two Competing Costs}

A single interaction event embedded in a causal diamond of spatial
edge $\lambda$ incurs two costs:

\begin{enumerate}
  \item \textbf{Bit cost $J_{\mathrm{bit}}$.} The information cost of
    the interaction event itself. In normalized cost accounting,
    \begin{equation}
      J_{\mathrm{bit}} = 1.
      \label{eq:Jbit-norm}
    \end{equation}
    The unnormalized form is $J_{\mathrm{bit}} = \ln\varphi \approx
    0.481$ (the natural information content of a $\varphi$-bit, as
    defined in~\eqref{eq:Jbit-def}). The normalization absorbs
    $\ln\varphi$ into the cost unit.

  \item \textbf{Curvature cost $J_{\mathrm{curv}}(\lambda)$.} The
    geometric cost from Section~\ref{sec:curvature-cost}:
    $J_{\mathrm{curv}}(\lambda) = 2\lambda^2$.
\end{enumerate}

\subsection{The Total Cost Functional}

\begin{equation}
  C(\lambda)
  = J_{\mathrm{bit}} + J_{\mathrm{curv}}(\lambda)
  = 1 + 2\lambda^2.
  \label{eq:total-cost}
\end{equation}

\subsection{The Minimal-Overhead Balance Condition}

The minimal-overhead principle requires the discrete
lattice to operate at the scale where neither cost channel carries
unnecessary surplus. The total cost $C(\lambda)$ sums a
scale-independent term (bit cost) and a scale-dependent term
(curvature cost, growing as $\lambda^2$). The equilibrium is at
the crossover:
\begin{equation}
  J_{\mathrm{bit}} = J_{\mathrm{curv}}(\lambda_0).
  \label{eq:balance}
\end{equation}
If $\lambda > \lambda_0$, the curvature cost exceeds the bit cost and
the lattice wastes geometric capacity. If $\lambda < \lambda_0$, the
bit cost exceeds the curvature cost and the lattice wastes information
capacity. The minimal-overhead principle excludes both regimes.

\begin{theorem}[Recognition length in dimensionless units]
\label{thm:lambda-dimless}
The unique positive solution to $J_{\mathrm{bit}} =
J_{\mathrm{curv}}(\lambda_0)$ is $\lambda_0 = 1/\sqrt{2}$.
\end{theorem}

\begin{proof}
$J_{\mathrm{curv}}(\lambda) = 2\lambda^2$ is strictly monotone
increasing for $\lambda > 0$, and $J_{\mathrm{bit}} = 1 > 0$. The
equation $2\lambda^2 = 1$ has exactly one positive solution:
$\lambda_0 = 1/\sqrt{2}$.
\end{proof}

\begin{remark}[No circular dependence]
This derivation uses only the bit cost (from $\varphi$-scaling) and
the curvature cost (from $Q_3$ geometry via Gauss-Bonnet). It does
not reference Newton's constant $G$, the Planck length $\ell_P$, or
any external physical input. $G$ will be \emph{defined} as a
consequence in Section~\ref{sec:consequences}.
\end{remark}

%% ============================================================
\section{Dimensional Restoration and the Planck Gate Identity}
\label{sec:dim-restore}

\subsection{From Dimensionless Balance to SI}

The balance condition~\eqref{eq:balance} is a statement in natural
units, where $\lambda_{\mathrm{rec}} = 1$, $c = 1$, and
$\hbar = \varphi^{-5}$. The balance point $\lambda_0 = 1/\sqrt{2}$
is dimensionless. To convert to SI, we must determine what $G$---and
hence the Planck length $\ell_P = \sqrt{\hbar G/c^3}$---equals.

The curvature-cost derivation \emph{determines} $G$. The factor of
$\pi$ relating $\lambda_{\mathrm{rec}}$ to $\ell_P$ comes from
Gauss-Bonnet.

\subsection{The Gauss-Bonnet Bridge}
\label{sec:GB-bridge}

The Gauss-Bonnet integral of $S^2$ is
$\int K\,dA = 4\pi$, carried by $|\kappa| = 4$ curvature modes
(Proposition~\ref{prop:kappa}). Each mode contributes a curvature
load
\begin{equation}
  \frac{\int K\,dA}{|\kappa|}
  = \frac{4\pi}{4} = \pi.
  \label{eq:GB-per-mode}
\end{equation}
The Planck area $\ell_P^2 = \hbar G/c^3$ sets the area of one
curvature quantum. Each mode carries $\pi$ such quanta (in natural
units where $\lambda_{\mathrm{rec}} = 1$); the area per quantum is
the inverse of the curvature load:
\begin{equation}
  \ell_P^2 = \frac{1}{\pi}
  \quad(\text{natural units}).
  \label{eq:ell-P-RS}
\end{equation}
Equivalently: the voxel face area ($\lambda_{\mathrm{rec}}^2 = 1$)
contains $\pi$ Planck cells per curvature mode.

\subsection{Newton's Constant as a Derived Quantity}
\label{sec:G-derived}

\begin{definition}[Newton's constant from $\lambda_{\mathrm{rec}}$]
\label{def:G-def}
From $\ell_P^2 = \hbar G/c^3$ and~\eqref{eq:ell-P-RS}:
\begin{equation}
  \boxed{
    G = \frac{c^3\, \lambda_{\mathrm{rec}}^2}{\pi\,\hbar}.
  }
  \label{eq:G-derived}
\end{equation}
In natural units ($\lambda_{\mathrm{rec}} = 1$, $c = 1$,
$\hbar = \varphi^{-5}$), this gives $G = \varphi^5/\pi$.
\end{definition}

\begin{corollary}[Einstein coupling constant]
\label{cor:kappa}
The Einstein gravitational coupling
$\kappa = 8\pi G/c^4$ evaluates to
\begin{equation}
  \boxed{\kappa = 8\varphi^5.}
  \label{eq:kappa}
\end{equation}
This is machine-verified as \texttt{kappa\_einstein\_eq} in
\texttt{Constants.lean} and \texttt{kappa\_rs\_closed\_form} in
\texttt{Gravity/ZeroParameterGravity.lean}
\cite{WashburnLean4Repo2026}.
\end{corollary}

\begin{proof}
$\kappa = 8\pi G/c^4 = 8\pi \cdot
\varphi^5/(\pi \cdot 1) = 8\varphi^5$
(using $G = \varphi^5/\pi$ and $c = 1$ in natural units).
\end{proof}

\subsection{The Recognition Length in SI}

Since $\ell_P^2 = \lambda_{\mathrm{rec}}^2/\pi$
(from~\eqref{eq:ell-P-RS}):
\begin{equation}
  \boxed{
    \lambda_{\mathrm{rec}}
    = \sqrt{\pi}\;\ell_P
    = \sqrt{\frac{\pi\,\hbar G}{c^3}}
    \approx 2.86 \times 10^{-35}\;\mathrm{m}.
  }
  \label{eq:lambda-rec}
\end{equation}

\subsection{The Planck Gate Identity}

\begin{equation}
  \boxed{
    \frac{c^3\,\lambda_{\mathrm{rec}}^2}{\hbar\,G}
    = \pi.
  }
  \label{eq:planck-gate}
\end{equation}

\begin{proof}[Verification]
Substituting $G = c^3 \lambda_{\mathrm{rec}}^2/(\pi\hbar)$:
$c^3\lambda_{\mathrm{rec}}^2/(\hbar \cdot c^3
\lambda_{\mathrm{rec}}^2/(\pi\hbar)) = \pi$.
This identity is machine-verified in Lean~4
(Appendix~\ref{app:lean}).
\end{proof}

%% ============================================================
\section{Consequences}
\label{sec:consequences}

\subsection{Position in the Forcing Chain}

With $\lambda_{\mathrm{rec}}$ and $G$ derived, the
chain~\eqref{eq:chain} continues to the fine-structure constant
\begin{equation}
  \alpha^{-1} = 4\pi \cdot 11 - w_8 \ln\varphi
    - \frac{103}{102\pi^5}
    \approx 137.035
    \quad\text{(Thomson limit)},
\end{equation}
and onward to the mass sector. Each step is algebraic once
$\lambda_{\mathrm{rec}}$ is in place.

\subsection{Quantum-Gravitational Bridge}

$\lambda_{\mathrm{rec}}$ sits at the balance point of the
curvature cost: the length scale where bit cost equals curvature cost.
Below $\lambda_{\mathrm{rec}}$, the $\hbar$ sector dominates. Above
$\lambda_{\mathrm{rec}}$, the $G$ sector dominates. The reciprocal
symmetry $J(x) = J(x^{-1})$ maps one regime to the other;
$\lambda_{\mathrm{rec}}$ is the fixed point.

%% ============================================================
\section{Independent Check: Bekenstein-Landauer Lower Bound}
\label{sec:bekenstein}

The Bekenstein-Hawking entropy bound $S \leq A/(4\ell_P^2)$ and the
Landauer principle jointly provide an \emph{independent}
information-theoretic constraint on any fundamental length scale,
derived without reference to the cost functional or Gauss-Bonnet
geometry. An earlier version of this argument appears in
\cite{Washburn2025IntrinsicLattice}.

\begin{axiomRS}[One-bit causal diamond]
$S_\diamond = k_B \ln 2$.
\end{axiomRS}

Saturating the Bekenstein-Hawking bound gives the minimum bounding area
capable of encoding one bit:
\begin{equation}
  A_\diamond \;=\; 4\,\ell_P^2\,\ln 2.
\end{equation}
The Bekenstein-Hawking formula is derived for spherical horizons, so the
natural bounding geometry is the sphere $S^2$ of radius
$\lambda_{\mathrm{BL}}$:
\begin{equation}
  4\pi\,\lambda_{\mathrm{BL}}^2 \;=\; 4\,\ell_P^2\,\ln 2
  \qquad\Longrightarrow\qquad
  \lambda_{\mathrm{BL}}
  \;=\; \sqrt{\frac{\ln 2}{\pi}}\;\ell_P
  \;\approx\; 7.6 \times 10^{-36}\;\mathrm{m}.
  \label{eq:lambda-BL}
\end{equation}

The curvature-extremum value is
$\lambda_{\mathrm{rec}} = \sqrt{\pi}\,\ell_P \approx 2.86 \times
10^{-35}$~m, a factor of
\begin{equation}
  \frac{\lambda_{\mathrm{rec}}}{\lambda_{\mathrm{BL}}}
  \;=\; \frac{\pi}{\sqrt{\ln 2}}
  \;\approx\; 3.77
\end{equation}
above the Bekenstein-Landauer floor. This factor has a direct
physical reading. The Bekenstein bound is an \emph{inequality}: it
sets the smallest sphere that \emph{can} encode one bit. The
cost-balance condition (Section~\ref{sec:main-result}) selects the
scale at which the lattice \emph{does} operate. That the operating
scale exceeds the information-theoretic minimum by a modest factor is
expected and healthy: it means every recognition voxel carries
comfortably more area than the Bekenstein minimum requires. The
Bekenstein-Landauer route therefore serves as a consistency floor, not
as a competing estimate of $\lambda_{\mathrm{rec}}$.

%% ============================================================
\section{Phenomenological Consistency Checks}
\label{sec:consistency}

Via the effective recognition length $\lambda_{\mathrm{eff}} =
\lambda_{\mathrm{micro}} \cdot f^{-1/4}$
\cite{Washburn2025DualDerivation}:

\begin{enumerate}
  \item \textbf{Stellar luminosity balance:}
    $\lambda_{\mathrm{eff}}^{(\star)} = 6.3 \times 10^{-5}$~m,
    $f^{(\star)} \simeq 3.2 \times 10^{-122}$.
  \item \textbf{Vacuum energy density:}
    $\lambda_{\mathrm{eff}}^{(\Lambda)} = 5.9 \times 10^{-5}$~m,
    $f^{(\Lambda)} \simeq 3.4 \times 10^{-122}$.
\end{enumerate}

The two routes agree within 7\%:
$f = (3.3 \pm 0.3) \times 10^{-122}$,
$\lambda_{\mathrm{eff}} \approx 60\;\mu\mathrm{m}$.
These are validations, not inputs.

%% ============================================================
\section{Discussion}
\label{sec:discussion}

\subsection{Summary of the Derivation}

Three ingredients, all internal to the derivation:
\begin{enumerate}
  \item Bit cost $J_{\mathrm{bit}} = 1$ (normalized), from
    $\varphi$-scaling (Theorem~\ref{thm:phi}).
  \item Curvature cost $J_{\mathrm{curv}} = 2\lambda^2$, from the
    $|\kappa| = 4$ curvature quanta on $Q_3$ and the Gauss-Bonnet
    normalization on $S^2$
    (Proposition~\ref{prop:kappa},
    Section~\ref{sec:curvature-cost}).
  \item Balance condition $J_{\mathrm{bit}} = J_{\mathrm{curv}}$
    from the minimal-overhead principle.
\end{enumerate}
The balance selects $\lambda_0 = 1/\sqrt{2}$ in natural units.
The Gauss-Bonnet bridge (Section~\ref{sec:GB-bridge}) identifies
$\ell_P^2 = 1/\pi$ in natural units, yielding
$\lambda_{\mathrm{rec}} = \sqrt{\pi}\,\ell_P$ and
$G = c^3 \lambda_{\mathrm{rec}}^2/(\pi\hbar)$, with the Einstein
coupling $\kappa = 8\varphi^5$.

\subsection{Why Dimensional Closure Was Hard}

$J(x)$ is scale-invariant: it depends only on ratios. The
curvature-extremum argument breaks this invariance by introducing a
cost quadratic in the physical edge length $\lambda$, not in a ratio.
When balanced against the fixed bit cost, it selects a unique scale.

\subsection{Scope and Assumptions}
\label{sec:scope}

The derivation depends on three classes of input, listed here with
their epistemic status:
\begin{enumerate}[label=(\alph*)]
  \item \textbf{Theorems.}  The uniqueness of $J(x)$
    (Theorem~\ref{thm:J-unique}) and the forcing of $\varphi$, $c$,
    $\hbar$ along the derivation chain~\eqref{eq:chain} are
    mathematical results, machine-verified in Lean~4
    \cite{WashburnLean4Repo2026}.
  \item \textbf{Framework results.}  The curvature token count
    $|\kappa| = 4$ (Proposition~\ref{prop:kappa}) and the
    Gauss-Bonnet normalization on $S^2$ follow from the $D = 3$ cube
    topology $Q_3$, but rely on a specific mode-decomposition and
    equal-weight argument that some readers may classify as modeling
    assumptions rather than unavoidable consequences.
  \item \textbf{Postulates.}  Two inputs are physical postulates, not
    theorems:
    \begin{itemize}
      \item The minimal-overhead balance condition
        $J_{\mathrm{bit}} = J_{\mathrm{curv}}$. This is the
        assertion that the lattice operates at the scale where
        informational and geometric costs are equal.
      \item The Planck-area bridge $\ell_P^2 = 1/\pi$
        (Section~\ref{sec:GB-bridge}), identifying the Planck area
        with the inverse of the Gauss-Bonnet curvature load per mode.
        This step connects the internal combinatorial geometry to the
        gravitational scale.
    \end{itemize}
\end{enumerate}
Given items (a)--(c), the output ($\lambda_{\mathrm{rec}}$, $G$,
$\kappa$) is unique with zero continuously adjustable parameters.
The paper does not claim that these postulates are themselves derived
from pure logic; it claims that, once accepted, they leave no
parameter free to tune.

\subsection{Curvature Scale versus Bekenstein Floor}

The curvature-extremum route gives the operating scale
$\lambda_{\mathrm{rec}} = \sqrt{\pi}\,\ell_P
\approx 2.86 \times 10^{-35}$~m.
The Bekenstein-Landauer route
(Section~\ref{sec:bekenstein}) gives an information-theoretic lower
bound $\lambda_{\mathrm{BL}} = \sqrt{\ln 2/\pi}\;\ell_P \approx 7.6
\times 10^{-36}$~m.  The ratio
$\lambda_{\mathrm{rec}}/\lambda_{\mathrm{BL}} =
\pi/\sqrt{\ln 2} \approx 3.77$ is not a discrepancy but a consistency
check: the Bekenstein bound is an inequality, and the lattice operating
above this floor by a moderate algebraic factor is the expected regime.
Both routes use spherical geometry ($S^2$) for the bounding surface, so
the comparison involves no ad hoc geometric choice.

The curvature-extremum value is the canonical result because it follows
from the cost functional and Gauss-Bonnet normalization alone, without
external bounds (Bekenstein-Hawking, Landauer), and it is consistent
with the machine-verified identity $\kappa = 8\varphi^5$ in the Lean
formalization \cite{WashburnLean4Repo2026}.

%% ============================================================
\section{Discussion and Conclusion}
\label{sec:conclusion}

With $\lambda_{\mathrm{rec}}$ derived, the framework determines both
dimensionless and dimensional constants. The number of independent
dimensional inputs drops from three ($G$, $\hbar$, $c$) to two
($\hbar$, $c$), with $G = c^3
\lambda_{\mathrm{rec}}^2/(\pi\hbar)$ and $\kappa = 8\varphi^5$.
The factor of $\pi$ originates in the Gauss-Bonnet curvature load
per mode on $S^2$.

$\lambda_{\mathrm{rec}}$ is where the informational cost of an
interaction event matches the geometric cost of the curvature it
creates. Below it, $\hbar$ dominates; above it, $G$ dominates.
The reciprocal symmetry $J(x) = J(x^{-1})$ maps one regime to the
other, and $\lambda_{\mathrm{rec}}$ is the fixed point. In the
broader RS program, this scale is the prerequisite for addressing
the mass spectrum: particle masses become expressible once a
dimensional anchor is available.

All prerequisite theorems and the balance-condition derivation are
machine-verified in Lean~4 with zero \texttt{sorry} statements
(Appendix~\ref{app:lean}).

%% ============================================================
\appendix

\section{Lean~4 Machine Verification}
\label{app:lean}

The full Lean~4 codebase is at
\url{https://github.com/jonwashburn/recognition-science}
\cite{WashburnLean4Repo2026}. It contains 122 Lean files across 10
directories, with 32 verification certificates bundling related
theorems. One axiom is used:
\texttt{aczel\_representation\_axiom} in
\texttt{Cost/FunctionalEquation.lean}, encoding the Acz\'el
smoothness condition for the d'Alembert equation (a standard result;
Acz\'el, 1966). All other results are proved unconditionally.

\subsection{What Lean Verifies}

The formalization covers:
\begin{enumerate}[label=(\roman*)]
  \item \textbf{Prerequisite theorems.} The uniqueness of $J$
    (Theorem~\ref{thm:J-unique}), the forcing of $\varphi$
    (Theorem~\ref{thm:phi}), and the forcing of $D = 3$
    (Theorem~\ref{thm:D3}), formalized with zero \texttt{sorry}
    statements.
  \item \textbf{Balance-condition derivation.} The cost balance
    $J_{\mathrm{bit}} = J_{\mathrm{curv}}(\lambda)$, i.e.,
    $1 = 2\lambda^2$, has the unique positive root
    $\lambda_0 = 1/\sqrt{2}$. This is formalized as
    \texttt{balance\_at\_lambda\_0} and
    \texttt{balance\_unique\_positive\_root} in
    \texttt{Constants/LambdaRecDerivation.lean}
    \cite{WashburnLean4Repo2026}. These proofs reference
    only the bit cost (= 1) and the curvature cost ($= 2\lambda^2$),
    \emph{without} reference to $G$, $\ell_P$, or any external
    constant.
  \item \textbf{Algebraic consistency.} \texttt{Constants.lean}
    defines $G := \lambda_{\mathrm{rec}}^2 c^3/(\pi\hbar)$ and
    verifies $\kappa = 8\varphi^5$ via
    \texttt{kappa\_einstein\_eq}.
    The Planck gate identity
    $c^3\lambda_{\mathrm{rec}}^2/(\hbar G) = \pi$
    is an algebraic consequence of this definition; its role is to
    confirm internal consistency.
\end{enumerate}

\subsection{What Lean Does Not Verify}

The physical postulates listed in Section~\ref{sec:scope}---the
balance condition $J_{\mathrm{bit}} = J_{\mathrm{curv}}$ and the
Planck-area bridge $\ell_P^2 = 1/\pi$---are \emph{inputs} to the
formalization, not outputs. Lean verifies the mathematical
consequences of these postulates (e.g., that $\lambda_0 = 1/\sqrt{2}$
is the unique root, that $\kappa = 8\varphi^5$ follows algebraically),
but it does not establish that the postulates are physically
inevitable. The question of whether the balance condition and the
Planck-area identification are unavoidable consequences of discrete
lattice dynamics, or whether they are substantive modeling choices,
remains open.

\subsection{Module Table}

Table~\ref{tab:lean} lists the modules from the public repository
\cite{WashburnLean4Repo2026} that formalize theorems cited in this
paper.

\begin{table}[ht]
\centering
\small
\begin{tabular}{@{}p{5.2cm}p{5.3cm}p{2cm}@{}}
\toprule
\textbf{Module path} & \textbf{Key theorems} & \textbf{sorry} \\
\midrule
\texttt{Cost/JcostCore.lean}
  & \texttt{Jcost\_symm},
    \texttt{Jcost\_nonneg},
    \texttt{Jcost\_pos\_of\_ne\_one},
    \texttt{Jcost\_strict\_mono\_\newline on\_one\_infty}
  & 0 \\[4pt]
\texttt{Cost/FunctionalEquation\newline .lean}
  & \texttt{washburn\_zlatanovic\_\newline uniqueness\_aczel},
    \texttt{dAlembert\_to\_ODE\_\newline theorem},
    \texttt{ode\_cosh\_uniqueness\_\newline contdiff}
  & 0 \\[4pt]
\texttt{Foundation/PhiForcing\newline .lean}
  & \texttt{phi\_forced},
    \texttt{golden\_constraint\_\newline unique}
  & 0 \\[4pt]
\texttt{Foundation/EightTick\newline .lean}
  & \texttt{phase\_eighth\_power\_\newline is\_one},
    \texttt{shift\_period\_8}
  & 0 \\[4pt]
\texttt{Foundation/Dimension\-Forcing.lean}
  & \texttt{eight\_tick\_forces\_D3},
    \texttt{dimension\_forced}
  & 0 \\[4pt]
\texttt{Constants.lean}
  & \texttt{phi\_sq\_eq},
    \texttt{hbar\_eq\_phi\_inv\_fifth},
    \texttt{G\_pos},
    \texttt{kappa\_einstein\_eq}
  & 0 \\[4pt]
\texttt{Constants/LambdaRec\-Derivation.lean}
  & \texttt{balance\_at\_lambda\_0},
    \texttt{balance\_unique\_\newline positive\_root},
    \texttt{lambda\_rec\_is\_forced}
  & 0 \\[4pt]
\texttt{Gravity/ZeroParameter\-Gravity.lean}
  & \texttt{kappa\_rs\_closed\_form},
    \texttt{kappa\_bounds},
    \texttt{equivalence\_principle\_\newline automatic}
  & 0 \\
\bottomrule
\end{tabular}
\caption{Lean~4 modules from
\texttt{github.com/jonwashburn/recognition-science}
that formalize theorems cited in this paper.
All paths are relative to \texttt{RecognitionScience/}.}
\label{tab:lean}
\end{table}

The repository also contains 32 verification certificates
(\texttt{Verification/} directory) that bundle related theorems.
Of particular relevance:
\texttt{NonCircularityCert} (verifies no experimental data enters
upstream),
\texttt{ZeroAdjustableParamsCert} (\texttt{knobsCount = 0}), and
\texttt{ForcingChainCert} (the complete derivation chain in one
structure).

\subsection{Correspondence Between Paper and Lean}

\begin{itemize}
  \item \textbf{Theorem~\ref{thm:J-unique}} (uniqueness of $J$):
    formalized as
    \texttt{washburn\_zlatanovic\_uniqueness\_aczel} in
    \texttt{Cost/FunctionalEquation.lean}. The proof uses the Acz\'el axiom
    (\texttt{aczel\_dAlembert\_smooth}) to establish smoothness, then
    derives the ODE $H'' = H$ via
    \texttt{dAlembert\_to\_ODE\_theorem}, and concludes $H = \cosh$
    via \texttt{ode\_cosh\_uniqueness\_contdiff}. The ODE uniqueness
    proof uses the integrating-factor diagonalization described in
    Section~\ref{sec:Jcost}: \texttt{deriv\_neg\_self\_zero} ($g' = -g
    \Rightarrow g = 0$) and \texttt{deriv\_pos\_self\_zero} ($h' = h
    \Rightarrow h = 0$).

  \item \textbf{Theorem~\ref{thm:phi}} (phi forcing): formalized as
    \texttt{phi\_forced} and \texttt{golden\_constraint\_unique} in
    \texttt{Foundation/PhiForcing.lean}. The proof uses
    $\varphi^2 = \varphi + 1$ (from \texttt{phi\_equation} in
    \texttt{Constants.lean}) and the algebraic identity
    $(r - \varphi)^2 \geq 0$ to force $r = \varphi$.

  \item \textbf{Theorem~\ref{thm:D3}} ($D = 3$ forcing): formalized
    as \texttt{eight\_tick\_forces\_D3} and \texttt{dimension\_forced}
    in \texttt{Foundation/DimensionForcing.lean}. The proof is by case
    analysis on $D$.

  \item \textbf{Theorem~\ref{thm:lambda-dimless}} (balance condition):
    formalized as \texttt{balance\_at\_lambda\_0} and
    \texttt{balance\_unique\_positive\_root} in
    \texttt{Constants/LambdaRecDerivation.lean}
    \cite{WashburnLean4Repo2026}. The module defines:
    \begin{align*}
      J_{\mathrm{bit}} &:= 1, \\
      J_{\mathrm{curv}}(\lambda) &:= 2\lambda^2, \\
      \text{balanceResidual}(\lambda) &:=
        J_{\mathrm{curv}}(\lambda) - J_{\mathrm{bit}}
        = 2\lambda^2 - 1.
    \end{align*}
    The proof establishes
    $\text{balanceResidual}(1/\sqrt{2}) = 0$ and uniqueness:
    for $\lambda > 0$,
    $\text{balanceResidual}(\lambda) = 0 \Leftrightarrow
    \lambda = 1/\sqrt{2}$.
    This is a \emph{non-circular} verification: it uses only the
    cost functionals, without reference to $G$ or $\ell_P$.

  \item \textbf{Natural-unit constants} (Definition~\ref{def:RS-constants}):
    all defined in \texttt{Constants.lean}
    \cite{WashburnLean4Repo2026}. Key definitions:
    \begin{align*}
      \texttt{phi} &:= (1 + \sqrt{5})/2, \\
      \texttt{hbar} &:= \texttt{cLagLock} \cdot \texttt{tau0}
        = \varphi^{-5} \cdot 1, \\
      \texttt{lambda\_rec} &:= \texttt{ell0} = 1, \\
      \texttt{G} &:= \texttt{lambda\_rec}^2
        \cdot \texttt{c}^3 / (\pi \cdot \texttt{hbar}).
    \end{align*}
    From these definitions,
    $\kappa = 8\pi G/c^4 = 8\varphi^5$ is proved as
    \texttt{kappa\_einstein\_eq}, and the Planck gate identity
    $c^3\lambda_{\mathrm{rec}}^2/(\hbar G) = \pi$
    follows algebraically.
\end{itemize}

%% ============================================================
\begin{thebibliography}{99}

\bibitem{WashburnZlatanovic2026Math935}
J.~Washburn and M.~Zlatanovi\'c,
\newblock Uniqueness of the Canonical Reciprocal Cost,
\newblock \emph{Mathematics} \textbf{14} (2026), 935.
\newblock \url{https://doi.org/10.3390/math14060935}

\bibitem{WashburnZlatanovicAllahyarov2026Inevitability}
J.~Washburn, M.~Zlatanovi\'c, and E.~Allahyarov,
\newblock The d'Alembert Inevitability Theorem,
\newblock preprint, March 11, 2026.
\newblock Placeholder reference for forthcoming arXiv or journal version.

\bibitem{WashburnBarghi2026Axioms151}
J.~Washburn and A.~R.~Barghi,
\newblock Reciprocal Convex Costs for Ratio Matching: Axiomatic
Characterization,
\newblock \emph{Axioms} \textbf{15} (2026), 151.
\newblock \url{https://doi.org/10.3390/axioms15020151}

\bibitem{WashburnZlatanovicAllahyarov2026RG}
J.~Washburn, M.~Zlatanovi\'c, and E.~Allahyarov,
\newblock Recognition Geometry,
\newblock \emph{Axioms} \textbf{15} (2026), 90.
\newblock \url{https://doi.org/10.3390/axioms15020090}

\bibitem{PardoGuerraEtAl2026}
S.~Pardo-Guerra, M.~Simons, A.~Thapa, and J.~Washburn,
\newblock Coherent Comparison as Information Cost: A Cost-First Ledger
Framework for Discrete Dynamics,
\newblock arXiv:2601.12194 (2026).

\bibitem{Washburn2025ParameterFree}
J.~Washburn,
\newblock Parameter-Free Derivations of Fundamental Dimensionless
Constants from Recognition Physics,
\newblock Recognition Physics Institute, February 2025.

\bibitem{Washburn2025IntrinsicLattice}
J.~Washburn,
\newblock The Recognition Length $\lambda_{\mathrm{rec}}$: Discovering
the Intrinsic Lattice Scale of Reality,
\newblock Recognition Physics Institute, 25 April 2025.

\bibitem{Washburn2025DualDerivation}
J.~Washburn,
\newblock Dual Derivation of the Effective Recognition Length
$\lambda_{\mathrm{eff}}$,
\newblock Recognition Physics Institute, May 2025.

\bibitem{Washburn2026HistoricalStudy}
J.~Washburn,
\newblock The Recognition Length in Recognition Science: Historical
Emergence, Dimensional Closure, and the Zero-Parameter Program,
\newblock Recognition Science Institute, March 2026.

\bibitem{WashburnLean4Repo2026}
J.~Washburn,
\newblock Recognition Science --- Lean~4 Formalization (source code),
\newblock \url{https://github.com/jonwashburn/recognition-science},
March 2026.
\newblock 122~Lean files, 32~verification certificates,
zero~\texttt{sorry}.

\end{thebibliography}

\end{document}
